% !PRINTED 472
\weberpage{472}

On the basis of the properties 1., 2., 3.\ we call the totality of the $\mu$ substitutions of the normal field $\Omega(\varrho)$ a \gesperrt{group}, and indeed a \gesperrt{finite group} of \gesperrt{degree} $\mu$. It is also called the \gesperrt{group of the substitutions of the field} $\Omega(\varrho)$.

With the general theory of groups we shall occupy ourselves in detail in later sections.

\begin{center}
\S.\ 148.
\end{center}

\begin{center}
\gesperrt{Permutation groups.}
\end{center}

In \S.\ 145 we have seen how, for an arbitrary equation in $\Omega$ of degree $m$
\weq{1}{F(x) = 0,}
of which it is only assumed that it has no equal roots, one can determine a normal field $\Omega(\varrho)$ to which all the roots of (1) belong. If, then, the roots of (1) are denoted by
\weq{2}{\alpha, \alpha_1, \alpha_2 \ldots \alpha_{m-1}}
then
\weq{3}{\alpha = \chi(\varrho), \quad \alpha_1 = \chi_1(\varrho), \ldots \alpha_{m-1} = \chi_{m-1}(\varrho),}
and the $\chi$ are rational functions in $\Omega$.

But if we set in place of $\varrho$ a quantity $\varrho_1$ conjugate to $\varrho$, that is, if we carry out a substitution $\sigma = (\varrho, \varrho_1)$, then every root $\alpha$ passes over into a definite other one, and never do two different ones pass into the same; for if
\[
F(\alpha) = F[\chi(\varrho)] = 0
\]
then by \S.\ 146, 1:
\[
F[\chi(\varrho_1)] = 0,
\]
that is, $\chi(\varrho_1)$ is also a root of the equation (1). But that two different quantities of the field $\Omega(\varrho)$ cannot pass over into the same quantity through a substitution $\sigma$ we have already shown in general in \S.\ 146.

\gesperrt{1. A substitution $\sigma$ therefore calls forth a certain permutation among the roots $\alpha$, in that each of the indices $0, 1, 2 \ldots m-1$ is replaced by a definite other one.}
